\section{Proof of the Bound for the Relative Zeta Value}\label{ce:section}

This section completes the proof of Theorem~\ref{an:ceiling}, the bound
$C_{\rm eff}$ for the relative zeta value that the proof of
Theorem~\ref{thm:main} uses. We apply Proposition~\ref{an:signed} with the
floors of Definition~\ref{dh:floors}, which form a system of floors by
Lemma~\ref{dh:floors-valid}, and with the kernel of
Definition~\ref{ce:kernel-def}. Lemma~\ref{ce:kernel} checks the conditions
on the kernel, and Lemma~\ref{ce:values} bounds the terms of the right-hand
side of Proposition~\ref{an:signed} other than $Z_{\rm sel}(1)$, from the census (Proposition~\ref{dh:census})
and the values of Propositions~\ref{an:genus-value} and~\ref{lv:values}.

\begin{definition}\label{ce:kernel-def}
Let $\sigma_1<\dots<\sigma_{21}$ and $c_1,\dots,c_{21}$ be the abscissae and
coefficients of Table~\ref{ce:kernel-table}; the $c_j$ are rational numbers,
given exactly in the witness file \texttt{witness\_0.043171.json}.
Let $b_{\rm TV}=8687735449/76582443134581$, so that
$b_{\rm TV}\approx1.1344291\cdot10^{-4}$, and let
$\mathsf K=\sum_jc_jg_{\sigma_j}$ be the kernel of
Proposition~\fref{an:signed} with $n=21$.
\end{definition}

\begin{table}[htbp]
\centering\footnotesize
\caption{The abscissae $\sigma_j$ and the coefficients $c_j$ of the kernel
$\mathsf K$, rounded to ten significant digits; the exact rational $c_j$ are
in the witness file. The coefficients were found by a linear programme over
$36$ abscissae, which left the others at $0$; the programme is only a
search, and Lemma~\ref{ce:kernel} certifies its output.}
\label{ce:kernel-table}
\begin{tabular}{rlr@{\qquad}rlr}
\toprule
$j$ & $\sigma_j$ & $c_j$ & $j$ & $\sigma_j$ & $c_j$\\
\midrule
$1$ & $1+1/43691$ & $0.02910294798$ & $12$ & $1+1/1499$ & $-151.8093424$\\
$2$ & $1+1/38178$ & $-0.08993884045$ & $13$ & $1+1/1000$ & $252.1350250$\\
$3$ & $1+1/29151$ & $0.2136909555$ & $14$ & $1+1/667$ & $-603.3009258$\\
$4$ & $1+1/22258$ & $-0.3750021755$ & $15$ & $1+1/510$ & $918.2080358$\\
$5$ & $1+1/14851$ & $0.8571627861$ & $16$ & $1+1/389$ & $-810.8699335$\\
$6$ & $1+1/9909$ & $-3.840291426$ & $17$ & $1+1/297$ & $442.3163431$\\
$7$ & $1+1/7566$ & $10.76698443$ & $18$ & $1+1/198$ & $-276.0370606$\\
$8$ & $1+1/5777$ & $-19.69796042$ & $19$ & $1+1/173$ & $175.9689941$\\
$9$ & $1+1/4411$ & $23.59901759$ & $20$ & $1+1/115$ & $-22.66409771$\\
$10$ & $1+1/2943$ & $-32.93912380$ & $21$ & $1+1/101$ & $8.602447576$\\
$11$ & $1+1/1964$ & $89.92828213$ & & &\\
\bottomrule
\end{tabular}
\end{table}

\begin{lemma}\label{ce:kernel}
The kernel $\mathsf K$ and the number $b_{\rm TV}$ of
Definition~\ref{ce:kernel-def} satisfy conditions (i)--(iii) of
Proposition~\ref{an:signed} for every real $q\ge q_{\min}=41^4$.
\end{lemma}

\begin{proof}[Proof (computer-assisted)]
Write $x=\log q$, and consider
\[
 \Phi_1=e^x\bigl(\mathsf K-g_1+b_{\rm TV}\,w\bigr)(e^x),\qquad
 \Phi_2=e^x\mathsf K(e^x),\qquad
 \Phi_3=e^x\frac{d}{dx}\mathsf K(e^x).
\]
Conditions (i)--(iii) say $\Phi_1\ge0$, $\Phi_2\ge0$ and $\Phi_3\le0$. The
supplementary program \texttt{signed.py} (function \texttt{certify})
evaluates them in Arb interval arithmetic in the forms
\[
 e^xg_s(e^x)=e^{-(s-1)x}\sum_{k\ge1}\frac{y^{k-1}}k,\qquad
 e^x\frac{d}{dx}g_s(e^x)=-s\,e^{-(s-1)x}\sum_{k\ge0}y^k,\qquad y=e^{-sx},
\]
which avoid cancellation, and
$e^xw(e^x)=2x\sum_{m\ge1}e^{-(m-1)x}/(1+e^{-mx})$. The series in $y$ are cut
after ten terms; since $x\ge\log q_{\min}>14$, the remainder and its first
three derivatives are less than $10^{-55}$, which is added as a radius. In
$w$ the terms with $m>6$ are dropped, which only lowers $\Phi_1$.

\emph{The range $\log q_{\min}\le x\le2\cdot10^6$.} This range is cut into
cells of width $1/20$ for $x<60$ and $x/100$ beyond. On a cell with midpoint
$x_0$ and half-width $r$, the program computes the Taylor coefficients of
orders $0$, $1$ and $2$ of $\Phi_k$ at $x_0$, and an enclosure of the third
derivative over the whole cell, by power series arithmetic over a ball; this
bounds $\Phi_k$ on the cell. A cell on which a bound fails is bisected. All
$1976$ cells pass, with $\Phi_1\ge1.86\cdot10^{-7}$,
$\Phi_2\ge3.73\cdot10^{-22}$ and $-\Phi_3\ge3.73\cdot10^{-22}$; the last two
minima occur near $x=2\cdot10^6$, where $\Phi_2$ is of order
$c_1e^{-(\sigma_1-1)x}$.

\emph{The range $x\ge X=2\cdot10^6$.} There $0\le y\le e^{-X}$, so
$1\le\sum_ky^{k-1}/k\le1+e^{-X}$ and $1\le\sum_ky^k\le1+2e^{-X}$. Since
$c_1>0$ and $\sigma_1$ is the smallest abscissa,
\begin{align*}
 e^{(\sigma_1-1)x}\Phi_2&\ge c_1-\sum_{j\ge2}|c_j|(1+e^{-X})e^{-(\sigma_j-\sigma_1)X},\\
 -e^{(\sigma_1-1)x}\Phi_3&\ge c_1\sigma_1-\sum_{j\ge2}|c_j|\sigma_j(1+2e^{-X})
 e^{-(\sigma_j-\sigma_1)X},
\end{align*}
and the program checks that both right sides are positive. Moreover
$e^xg_1(e^x)\le1+e^{-X}$ and $e^xw(e^x)\ge2x(1-e^{-X})$, and
$2b_{\rm TV}X>1+3e^{-X}$, so $g_1-b_{\rm TV}w<0\le\mathsf K$ there, which is
(i).
\end{proof}

\begin{lemma}\label{ce:values}
For the floors of Definition~\ref{dh:floors}, with floor sum $Y_{\rm fl}$,
and the kernel of Definition~\ref{ce:kernel-def},
\[
 \sum_{j=1}^{21}c_jY_{\rm fl}(\sigma_j)+b_{\rm TV}\,(2\kappa_\infty-T_{\rm sel})
 <0.000607837.
\]
\end{lemma}

\begin{proof}[Proof (computer-assisted)]
For $N>1$ put $T_{\mathbf c}(N)=\sum_jc_jT_{\sigma_j}(N)$, with $T_\sigma$ as
in Lemma~\ref{dh:floor-sum}. By that lemma,
\[
 \sum_jc_jY_{\rm fl}(\sigma_j)=\sum_jc_jY_{\rm cor}(\sigma_j)
 -\sum_{\mathfrak p\in\mathcal P_2,\ \mathrm N\mathfrak p>1.2\cdot10^7}
 T_{\mathbf c}(\mathrm N\mathfrak p),
\]
where $Y_{\rm cor}(\sigma)=Y_{W''}(\sigma)-\Delta_{\rm sel}(\sigma)
-Z_{\rm sel}(\sigma)-\Delta_4(\sigma)-\sum T_\sigma(\mathrm N\mathfrak p)$, the
last sum over the $136$ primes of $\mathcal P_2$ of norm at most
$1.2\cdot10^7$. The program \texttt{signed3.py} (with \texttt{signed.py})
computes $Y_{\rm cor}(\sigma_j)$ by \texttt{dihedral\_ceiling3.py} from the
enclosures of Propositions~\ref{lv:values} and~\ref{an:genus-value}, with
$\Delta_{\rm sel}$ given by Lemma~\ref{dh:delta-sel}, prime by prime for the
corrections, in Arb; since the $c_j$ have both signs, each
$Y_{\rm cor}(\sigma_j)$ is used as a two-sided enclosure. This gives
$\sum_jc_jY_{\rm cor}(\sigma_j)\le0.000561519780954$ (the enclosure is
$[0.000561490206974,\,0.000561519780954]$; only the upper end is used).

The primes of $\mathcal P_2$ above $1.2\cdot10^7$ are known through their bin
counts (Proposition~\ref{dh:census}). The program groups the nonempty bins
into runs of at most eight consecutive bins. On a run, every counted norm
lies in an interval $[{\rm lo},{\rm hi}]$, which includes one bin of slack on
each side, with midpoint $N_0$ and half-length $r$, and by the mean value
theorem
\[
 T_{\mathbf c}(N)\in T_{\mathbf c}(N_0)+T_{\mathbf c}'([{\rm lo},{\rm hi}])\,[-r,r],
\]
which is multiplied by the exact count of the run. Evaluating the large
signed coefficients only in the derivative term keeps the enclosure narrow.
This gives $\sum T_{\mathbf c}(\mathrm N\mathfrak p)\ge0.0000757399496680$
(the enclosure is $[0.0000757399496680,\,0.0000757536213684]$; only the lower
end is used). Finally $2\kappa_\infty-T_{\rm sel}=1.07342687156367\ldots$ is
computed in Arb from the definitions, so that
$b_{\rm TV}(2\kappa_\infty-T_{\rm sel})\le0.000121772671416$, and
\begin{align*}
 &0.000561519780954-0.0000757399496680+0.000121772671416\\
 &\qquad=0.000607552502702<0.000607837.
\end{align*}
The displayed endpoints are rounded outward from the Arb enclosures; the
supplementary program \texttt{ce\_values\_digits.py} prints them with all
digits, and the upper end of the total computed in Arb is
$0.00060755250270059\ldots$. The bound of the lemma is stated with a margin
of about $2.8\cdot10^{-7}$, so that it is not affected by changes of the
enclosures at the level of the rounding errors that their computation
accounts for.
\end{proof}

\begin{proof}[Proof of Theorem~\ref{an:ceiling}]
By Lemmas~\ref{dh:floors-valid} and~\ref{ce:kernel}, Proposition~\ref{an:signed}
applies to the floors of Definition~\ref{dh:floors} and the kernel of
Definition~\ref{ce:kernel-def}. By Lemmas~\ref{an:selected-lemma}
and~\ref{ce:values},
\begin{align*}
 Z_{\rm sel}(1)+\sum_jc_jY_{\rm fl}(\sigma_j)+b_{\rm TV}\,(2\kappa_\infty-T_{\rm sel})
 &<0.0416684840321405+0.000607837\\
 &<0.0422763211<\Ceff.
\end{align*}
The interval evaluation of the same expression from the unrounded
enclosures, by \texttt{signed.py}, gives the upper bound
$0.04227603653484\ldots$. Proposition~\ref{an:signed} with $C'=C_{\rm eff}$
proves the theorem.
\end{proof}

\begin{remark}\label{ce:comparison}
For comparison, the comparison field $\EW$ with a single abscissa does less.
At $\sigma=1001/1000$, the bound of Proposition~\ref{an:prime-budget} with
the termwise bound $Y_{W''}-\Delta_{\rm sel}-\Delta_4$, the census credit
$\sum_{\mathfrak p\in\mathcal P_2}T_\sigma(\mathrm N\mathfrak p)$ of
Lemma~\ref{dh:floor-sum}, and the adaptive credits of
Definition~\ref{lim:credits}, as computed by \texttt{dihedral\_ceiling3.py},
is $0.0425629$; the signed kernel lowers this by $2.9\cdot10^{-4}$. The
linear programme that chose the kernel is a floating-point search over $36$
abscissae; the proof uses only its rounded output, checked by
Lemma~\fref{ce:kernel}.
\end{remark}
